\chapter{Answers to the exercises}
\label{ch:svar}

Here are the answers to all the exercises. Read them after you have tried for yourselves --- an
answer you have read is not the same thing as an answer you have arrived at, and that difference is
what the exercises are for.

Where an answer rests on a calculation, the calculation is included, not just the result.

\section{The bus}

\solution{ex:1:1}{Wiring}
\ans{a} $8 \cdot 7 / 2 = 28$ connections.
\ans{b} 8 attachments, one per unit.
\ans{c} We look for $n$ where $n(n-1)/2 = 2n$, so $n - 1 = 4$ and $n = 5$. Already at five units
point-to-point is twice as much copper, and the gap then grows quickly.

\solution{ex:1:2}{The bus level}
\ans{a} \dom. The bus is dominant as soon as at least one node pulls it down, and node B does.
\ans{b} Node A and node C. Both sent recessive and read back dominant. Node D sent nothing and has
nothing to compare with.
\ans{c} That someone else is sending at the same time, and that the other one has priority. If it
happens in the arbitration field, A and C have lost arbitration and become receivers.

\solution{ex:1:3}{Termination}
\ans{a} There are three termination resistors instead of two: three 120~\textohm{} in parallel give
40~\textohm. Most likely a node in the middle of the bus has a resistor switched in.
\ans{b} Too low a termination attenuates the signal but does not destroy it on a short cable at
moderate speed --- the edges still have time to settle. The fault shows when the cable gets longer,
the speed higher or the environment noisier, that is, typically when the design moves off the lab
bench.

\solution{ex:1:4}{The asymmetry}
\ans{a} Yes. Arbitration only requires that \emph{one} of the two states wins unambiguously; which
one does not matter to the mechanism.
\ans{b} The highest identifier would win, since a one would then overwrite a zero.
\ans{c} Think about what the bus does when nobody is sending: it rests in the state nobody drives,
that is, recessive. With the chosen order the idle state is the \emph{losing} state, which means a
node that stops sending automatically withdraws. Were it the other way round, a silent bus would
look like a node winning everything, and a node that lost power would block the bus.

\section{The frame}

\solution{ex:2:1}{Read the frame}
\ans{a} $1 + 12 + 6 + 40 + 16 + 2 + 7 = 84$ bits. SOF 1, arbitration field 12 (11 identifier +
RTR), control field 6, data $5 \cdot 8 = 40$, CRC field 16 (15 + delimiter), ACK field 2, EOF 7.
\ans{b} 84~µs at 1~Mbit/s, before stuffing. With stuff bits and the three bits of intermission it is
somewhat more in practice.
\ans{c} \code{0x2A6}, since the lower numerical value wins.

\solution{ex:2:2}{No address}
\ans{a} One. The sensor sends the temperature once, with an identifier that describes the
\emph{content}.
\ans{b} Two frames, one to each receiver, or a broadcast address that the protocol has to define
specially.
\ans{c} Nothing. The third node simply starts listening for the same identifier. That is CAN's
strongest argument: receivers can be added without the sender knowing about it.

\solution{ex:2:3}{The acknowledgement}
\ans{a} That at least one node received the frame with a correct CRC.
\ans{b} How many did, which ones they were, whether any of them cared about the content, and
whether anyone actually processed it. The acknowledgement concerns the transfer, not the
processing.
\ans{c} It reads recessive, since there is nobody else to pull the slot down. A standard-compliant
controller treats that as an acknowledgement error, increases its transmit error counter and resends
--- which is why a lone node on a CAN bus sends the same frame over and over until someone answers.

\solution{ex:2:4}{DLC}
\ans{a} The field sits in the control field, whose width was fixed, and four bits were what was
available. That only 0--8 are used is the price of a field that is even in bit count rather than
stingy.
\ans{b} The driver's. The register layer passes the value on as it is, and the controller sequences
according to it --- neither of them validates. Putting the check in both places is exactly how the
two copies of the rule drift apart; it belongs where the frame is created.

\section{Bit stuffing and CRC-15}

\solution{ex:3:1}{Stuff a stream}
\ans{a} \code{0000011111} becomes \code{000001111101}: a stuff bit after the five zeros and one
after the five ones. 10 bits become 12.
\ans{b} \code{1111111111} becomes \code{111110111110}: a stuff bit after the first five ones, and
since the stuff bit itself breaks the run, the following five ones are counted from the start again,
which gives one more. 10 bits become 12.
\ans{c} \code{0101010101} is unchanged. No run reaches five, so no stuff bit is needed. 10 bits
remain 10.

\solution{ex:3:2}{Destuff}
\ans{a} The bit at position 6 (the zero after five ones) and the bit at position 11 (the one after
five zeros). Positions are counted from one.
\ans{b} \code{1111100000}.
\ans{c} A stuff error, after the sixth bit. After five equal bits the next one \emph{must} be the
opposite; if it is not, the bit stream is broken, and the receiver rejects the frame.

\solution{ex:3:3}{The CRC engine}
The register after each bit, starting from zero:

\begin{narrowtable}{@{}llll@{}}
\thead{Bit} & \thead{Value} & \thead{Register after} & \thead{Hex} \\ \midrule
1 & 1 & \code{100010110011001} & \code{0x4599} \\
2 & 1 & \code{000101100110010} & \code{0x0B32} \\
3 & 0 & \code{001011001100100} & \code{0x1664} \\
4 & 1 & \code{110100101010001} & \code{0x6951} \\
\end{narrowtable}

Notice steps 2 and 3: when the register's top bit and the incoming bit are equal, the feedback is
zero, and the step is a plain shift with no XOR.

\solution{ex:3:4}{Why zero?}
\ans{a} The checksum is the remainder of division by the polynomial. Appending the remainder to the
dividend makes it divide evenly, and an even division leaves a remainder of zero. The receiver
therefore needs no comparison --- it just checks whether the register is cleared.
\ans{b} The same module can be used on both the transmit and receive sides, with no mode switch and
no separate comparator circuit. In hardware that is a tangible saving, and one source of errors
fewer.
\ans{c} The transmitter's and receiver's checksums are then computed over different bit streams,
since the stuff bits depend on the data. The result is that every frame containing at least one
stuff bit is rejected as a CRC error --- that is, almost all of them. The symptom is confusing:
short frames with alternating data work, longer frames do not.

\section{Arbitration}

\solution{ex:4:1}{Two nodes}
\ans{a} \code{0x155} is \code{00101010101} and \code{0x154} is \code{00101010100}.
\ans{b} At bit 11, the last bit of the identifier.
\ans{c} B wins, since it sends dominant where A sends recessive. A detects that immediately, stops
sending, becomes a receiver for B's frame and tries again when the bus has become idle.

\solution{ex:4:2}{Three nodes}
In binary: \code{0x400} is \code{10000000000}, \code{0x200} is \code{01000000000} and \code{0x201} is
\code{01000000001}.
\ans{a} At bit 1. A sends recessive while B and C send dominant, so A drops out straight away.
\ans{b} At bit 11. B and C are identical for ten bits and differ only in the last one.
\ans{c} B, with the identifier \code{0x200}.

\solution{ex:4:3}{The same identifier}
\ans{a} Nothing. The bits are identical all through the arbitration field, so neither node sees any
difference between itself and the bus. Both believe they won.
\ans{b} As soon as the data fields differ, one node sends recessive and reads dominant. Outside the
arbitration field that is not a lost arbitration but a \emph{bit error}: the node aborts with an
error frame, both frames are destroyed, and both nodes resend --- whereupon the same thing happens
again.
\ans{c} Because it cannot arise in operation of a correctly configured bus. Identifiers are assigned
per sending node during system design; two nodes with the same identifier are a mistake in the
assignment, not something that happens by itself.

\solution{ex:4:4}{Starvation}
\ans{a} It loses every arbitration against every other node, since \code{0x7FF} is the highest
possible value and so the lowest priority. If the bus is always busy, it never gets through.
\ans{b} No. The protocol guarantees that the \emph{highest}-priority frame gets through within a
predictable time, and that is what determinism means in CAN. That low-priority traffic can starve is
a consequence of the priority order, and that order is the point.
\ans{c} For example: keep the bus load down (30--50~\% is the rule of thumb in vehicle systems);
plan the identifiers by what is time-critical; make transmissions periodic so that no node sends as
often as it can; or split the traffic across several buses.

\section{Bit timing}

\solution{ex:5:1}{Work out the timing}
\ans{a} $50\,000\,000 / 1\,000\,000 = 50$ clock cycles per bit.
\ans{b} $50\,000\,000 / 500\,000 = 100$ clock cycles per bit.
\ans{c} $50 \cdot 0.7 = 35$, that is, at tick 35 of 50.

\solution{ex:5:2}{The sample point}
\ans{a} So that the signal has time to reach the farthest node on the bus and come back before the
bit is read. That is the propagation segment's job, and the sample point has to come after it.
\ans{b} During arbitration a node would read the bus before another node's dominant bit had arrived,
and so believe it had won when it had lost. Two nodes would then send at the same time for the rest
of the frame.
\ans{c} The propagation segment. It is the segment that represents the travel time, and that is why
cable length and bit rate are linked.

\solution{ex:5:3}{Stuffing and sync}
\ans{a} Five.
\ans{b} Bit stuffing, the rule of five. After five equal bits a bit of the opposite value is always
inserted, which guarantees an edge.
\ans{c} The receivers could go long stretches without edges to resynchronise on, and drift until
they sample the wrong bit. In practice that would require a clock line or considerably more accurate
oscillators --- which is the very reason stuffing exists.

\solution{ex:5:4}{Two symptoms}
\ans{a} Nothing works. The nodes are out of phase from SOF onwards, so every frame is gibberish at
the other end. The fault shows immediately and is therefore the kinder of the two.
\ans{b} Most bits are read the same, but those close to an edge or affected by propagation delay are
read differently. The symptom is sporadic CRC and stuff errors that get worse with a longer cable and
higher load. It is harder to debug because it looks like a bad cable or a noisy environment, and
because it works well enough to make you believe the basic setup is right.

\section{Error handling}

\solution{ex:6:1}{Who sees the error?}
\ans{a} CRC error. The receivers detect it; the transmitter learns of it indirectly, through nobody
acknowledging.
\ans{b} Form error. Every node sees it, since every node knows which bits must be recessive.
\ans{c} Acknowledgement error. Only the transmitter sees it: it lets go of the ACK slot and reads back
recessive.
\ans{d} Stuff error. Every node sees it, since every node counts equal bits in a row.

\solution{ex:6:2}{The counters}
\ans{a} $8 \cdot 8 = 64$.
\ans{b} $64 - 20 = 44$.
\ans{c} Error active in both cases, since the limit for error passive is at 127.

Notice the proportions: eight errors cost 64, and it takes 64 successful transmissions to work them
off. The counter is deliberately slower on the way down than on the way up.

\solution{ex:6:3}{The asymmetry}
\ans{a} Because the counters do not measure how many errors have occurred, but how likely it is that
the error is the node's own. A transmitting node is involved in every error it sees; a receiving
node is one of many that saw the same error. The one sending is therefore a more likely source of
the error, and is penalised accordingly.
\ans{b} The broken node is involved in \emph{every} error and climbs eight steps at a time, while the
healthy node on a noisy bus rises one step at a time and has time to work it off between
disturbances. After a while the broken node is error passive or bus off while the healthy one is
still error active.

\solution{ex:6:4}{Bus off}
\ans{a} Because a node with a broken transmitter reports errors on every frame --- including the
frames that are entirely correct. Its error frames are dominant and destroy other nodes' traffic, so
a single broken node can take down the whole bus.
\ans{b} It would have become unusable. Every frame from every node would have been aborted by the
broken node's error flag, and all traffic would have stopped.
\ans{c} Because the fault may be transient: a loose connector, a disturbance, a voltage dip. A
permanent shutdown would require someone to go out and restart the node by hand, even when the fault
is long gone.

\section{The simplified controller}

\solution{ex:7:1}{What is missing}
\ans{a} Complete: the loser becomes a receiver and receives the winning frame, and resends its own
automatically when the bus has become idle. Simplified: the loser sets its error flag, leaves the
transmitter role and goes idle --- it does not receive the frame that won, and does not resend.
\ans{b} Complete: the receiver aborts with an error frame, every node discards the frame, and the
transmitter resends. Simplified: the receiver rejects the frame and sets the error flag. No error
frame is sent, so the transmitter knows nothing.
\ans{c} Complete: the transmitter detects the acknowledgement error and resends. Simplified: the
transmitter never reads back the ACK slot, so the transfer completes just as if all had gone well.
\ans{d} Complete: both end up in the receive queue. Simplified: the second overwrites the first if the
first has not yet been acknowledged, and the newest wins.

\solution{ex:7:2}{Sticky bits}
\ans{a} $40 \cdot 10^{-6} \cdot 50 \cdot 10^{6} = 2000$ clock cycles.
\ans{b} Practically nothing. The pulse is one cycle long, and the probability that a poll happens to
hit exactly that cycle is about 1 in 2000. The software would miss almost every event, and the few
it caught would look random.
\ans{c} Because an automatic clear on read makes it impossible to read the status twice, and so
impossible to write code that first checks a bit and then acts on it. An explicit clear also
separates ``I have seen the event'' from ``I have handled it'' --- and it is the software, not the
hardware, that knows when the latter is true.

\solution{ex:7:3}{Two steps}
\ans{a} It confuses \emph{the port is free again} with \emph{the frame got through}. The status bit
comes back both after a successful transmission and after a lost arbitration.
\ans{b} Wait for the port to become free, and \emph{then} read the error flag. Only that combination
says how it went.
\ans{c} Lost arbitration, stuff error, failed CRC check, and a received remote frame. All four give
the same single bit, so the driver cannot tell them apart.

\solution{ex:7:4}{The boundary}
\ans{a} For example by buffering received frames in a queue, so that the software never noticed that
the controller holds only one. Or by automatically clearing the status bit on read, so that the
software did not have to acknowledge.
\ans{b} Hidden behaviour is behaviour you cannot reason about. A queue that sometimes overflows gives
frames that vanish without a trace, and the debugging starts in the wrong place: in the software,
which did the right thing. A limitation you know about can be planned for; one you do not know about
can only be run into.
\ans{c} That an interface should be \emph{honest} rather than helpful. A layer that hides the
properties of the underlying hardware does not move the problem, it only makes it invisible --- and
what is invisible cannot be tested, measured or reasoned about. The same principle applies to a
driver API: it should say what the hardware does, not what you wish it did.
